\begin{table*}[t]
  \centering\small
  \begin{tabularx}{\textwidth}{@{}>{\raggedright\arraybackslash}Xrrrrr@{}}
    \toprule
    Excess size (\%) & s0 & s1 & s2 & Pooled & sd \\
    \midrule
    no prompt comments & $+59.9$ & $+49.8$ & $+71.6$ & $+60.4$ & 10.9 \\
    comment-shaped noise & $+54.9$ & $+46.0$ & $+58.8$ & $+53.2$ & 6.5 \\
    informative comments & $-0.1$ & $-2.1$ & $-15.1$ & $-5.8$ & 8.2 \\
    \midrule
    Satisfaction (\%) & s0 & s1 & s2 & Pooled & sd \\
    \midrule
    no prompt comments & $38.8$ & $38.3$ & $40.0$ & $39.0$ & 0.9 \\
    comment-shaped noise & $40.0$ & $40.0$ & $40.8$ & $40.3$ & 0.4 \\
    informative comments & $38.9$ & $36.8$ & $39.0$ & $38.2$ & 1.3 \\
    \midrule
    Criterion of record & s0 & s1 & s2 & Pooled & \\
    \midrule
    (i) uncommented arm ratchets & $\checkmark$ & $\checkmark$ & $\checkmark$ & $\checkmark$ & \\
    (iv) noise $\approx$ uncommented & $\checkmark$ & $\checkmark$ & $\checkmark$ & $\checkmark$ & \\
    \bottomrule
  \end{tabularx}
  \caption{
    Every criterion passes on every world draw. The upper block gives final excess size
    $\frac{|D_T|}{|D_\star|}{-}1$ per arm, with one column per seed, then the pooled value and the
    across-seed standard deviation in pp. Each seed is a fresh permutation sample of IFEval
    items, so the seeds are independent draws of the world pool and not repeated runs on one pool.
    The middle block gives last-three-round constraint satisfaction as a level, so the size result
    is readable beside the correctness it was bought at. The lower block gives the criteria of
    record. We fixed these before the run and evaluate
    each one per seed, since a criterion that passed pooled but failed on one seed of three would
    be an averaging artifact, and a gap here would show it. Across seeds, the uncommented arm varies little
    against the effect this design is powered on, which is why we read the pooled numbers in
    \S\ref{sec:protocol}.
  }
  \label{tab:seeds}
\end{table*}
